\section{Метод решения}

Логика программы разделена на независимые модули:
\begin{enumerate}
  \item \emph{вычислительное ядро} \texttt{src/matrix\_filter.h/.c} --- медиана одного окна,
    фильтрация пикселя, блока строк и всей матрицы; сортировка окна элементов --- быстрая
    сортировка \texttt{src/qsort.h/.c};
  \item \emph{ввод/вывод} \texttt{src/matrix\_io.h/.c} --- матрица из файла, с клавиатуры или
    генерация; \texttt{src/util.h/.c} --- \texttt{die()} и строгий разбор аргументов;
  \item \emph{пул потоков} \texttt{src/pool.h/.c} --- распределение блоков между потоками,
    синхронизация итераций и измерение времени;
  \item \emph{измерения} \texttt{src/bench.h/.c} --- таймер и статистика повторов
    (минимум, медиана);
  \item \emph{прикладной уровень} \texttt{src/app/main.c} --- разбор ключей, одиночный прогон,
    режим \texttt{--bench}, проверка результата.
\end{enumerate}

Алгоритм параллельного прогона (\texttt{filter\_run}):

\begin{enumerate}
  \item матрица режется на непересекающиеся блоки из целых строк; по умолчанию
    \texttt{block\_rows} = \texttt{rows / (4P)}, так что число блоков $\approx 4P$ и работа
    равномерно дробится между потоками;
  \item создаётся не более $P-1$ рабочих потоков: лимит \texttt{--threads P} действует вместе
    с главным потоком; при $P=1$ считает сам \texttt{main} без барьеров и мьютексов;
  \item все потоки, включая главный, работают по одному циклу: захват номера блока из общей
    очереди под мьютексом, применение медианного фильтра к строкам блока \emph{без} мьютекса
    (блоки не пересекаются --- записей в общие данные нет), затем барьер итерации и
    переключение локальных буферов \texttt{a} $\leftrightarrow$ \texttt{b};
  \item барьер \texttt{pthread\_barrier\_t} на \texttt{P} участников --- ровно один на
    итерацию: пока все блоки текущей итерации не обработаны, никто не читает результат
    предыдущей. Двойной буфер нужен, чтобы фильтрация шла по данным предыдущей итерации, а не
    уже перекрашенным текущей; исходная матрица \texttt{src} не изменяется;
  \item счётчик взятых блоков \texttt{next\_block} не сбрасывается, а растёт сквозным: на
    итерации \texttt{iter} доступны блоки с номерами
    $[\texttt{iter}\cdot T,\ (\texttt{iter}+1)\cdot T)$, где $T$ --- общее число блоков.
    Поэтому очередь одна на весь прогон, а критическая секция --- только захват номера блока;
  \item время измеряется \texttt{clock\_gettime(CLOCK\_MONOTONIC)} от начала создания потоков
    до конца последней итерации: \texttt{join} и копирование результата в измеряемый интервал
    не входят, создание потоков --- входит.
\end{enumerate}

Общий ресурс, к которому обращаются несколько потоков, --- очередь блоков
(\texttt{next\_block} под \texttt{pthread\_mutex}). Итоговые данные детерминированы: обработка
блока не зависит от разбиения и числа потоков, поэтому параллельный результат можно сверить с
последовательным эталоном \texttt{filter\_run\_serial} (ключ \texttt{--check}, вне измеряемого
интервала).

\section{Описание программы}

\textbf{Разделение по файлам:}

\begin{center}
\begin{tabular}{ll}
\texttt{src/matrix\_filter.h/.c} & медианное окно, фильтр блока строк и всей матрицы \\
\texttt{src/qsort.h/.c}          & быстрая сортировка окна \\
\texttt{src/matrix\_io.h/.c}     & ввод из файла / с клавиатуры / генерация, вывод в файл \\
\texttt{src/util.h/.c}           & \texttt{die()}, разбор аргументов командной строки \\
\texttt{src/pool.h/.c}           & пул потоков: очередь, мьютекс, барьеры, таймер \\
\texttt{src/bench.h/.c}          & таймер и статистика повторов \\
\texttt{src/app/main.c}          & ключи запуска, одиночный прогон, \texttt{--bench}, \texttt{--check} \\
\end{tabular}
\end{center}

\textbf{Основной тип данных} --- \texttt{filter\_job}: исходная и результирующая матрицы,
двойной рабочий буфер \texttt{a}/\texttt{b}, размеры \texttt{rows/cols/window/iters},
разбиение \texttt{block\_rows}/\texttt{total\_blocks}, мьютекс и барьер, сквозной счётчик
\texttt{next\_block}, измеренное время \texttt{elapsed}.

\textbf{Функции пула потоков и их семантика:}

\begin{center}
\begin{tabular}{ll}
\texttt{filter\_job\_init()}        & инициализация задачи, создание мьютекса и буферов \\
\texttt{filter\_job\_set\_block\_rows()} & установка размера блока и пересчёт числа блоков \\
\texttt{filter\_run()}              & параллельный прогон на \texttt{threads} потоках \\
\texttt{filter\_run\_serial()}      & последовательный эталон \\
\texttt{bench\_stats()}             & минимум и медиана по повторам \\
\end{tabular}
\end{center}

\textbf{Используемые примитивы POSIX:}
\texttt{pthread\_create}, \texttt{pthread\_join},
\texttt{pthread\_mutex\_init/lock/unlock/destroy},
\texttt{pthread\_barrier\_init/wait/destroy},
\texttt{clock\_gettime(CLOCK\_MONOTONIC)}.

\textbf{Обработка ошибок:} функции ввода возвращают \texttt{void} --- при любой ошибке
(файл не открылся, чисел не хватило, вместо числа мусор) печатается сообщение и вызывается
\texttt{exit(EXIT\_FAILURE)}; неудача \texttt{pthread\_create}/\texttt{pthread\_barrier\_init}
также приводит к аварийному завершению через \texttt{die()}. Ключи проверяются строго
(нечётность окна, ровно один источник данных, несовместимые комбинации).
